\documentclass[11pt,a4paper]{article}
\usepackage[margin=2.4cm]{geometry}
\usepackage{graphicx}
\usepackage{booktabs}
\usepackage{amsmath}
\usepackage[hidelinks]{hyperref}
\usepackage{microtype}

\title{Module E: a recombination-matched haplodiploid neutral null for
ancestry sorting in replicate \textit{Formica} hybrid populations}
\author{}
\date{}

\begin{document}
\maketitle
\vspace{-2.2em}

\begin{abstract}
\noindent
Modules A--C established that parental ancestry sorts in a directional,
partly predictable way across 20 replicate \textit{F.\ polyctena}\,$\times$\,\textit{F.\
aquilonia} hybrid populations, but those results are descriptive: without a
neutral null they cannot distinguish selection from drift. Module E builds that
null. Founders are seeded from \emph{real phased parental haplotypes} rather than
independent allele-frequency draws, so within-species linkage disequilibrium is
preserved; recombination follows the empirical genetic map; the model is
haplodiploid; and all selection is removed, leaving drift and recombination as
the only forces. The key design element is an \emph{internal control}: markers are
stratified by DiagnosticIndex (DI), so low-DI markers---which cannot be driven by
ancestry sorting---calibrate the demography within the same runs that test the
high-DI markers. Observed among-population differentiation rises steeply with DI
($F_{ST}=0.051\rightarrow0.315$) while the neutral null is essentially flat
($0.191\rightarrow0.156$), and the fraction of sorted loci exceeds the null by
$\sim$1.5-fold at low DI rising to $>$100-fold at high DI. Because drift is
DI-blind by construction, a DI-scaled excess cannot be produced by demographic
misspecification. The null does not yet reproduce the neutral background itself
(it over-differentiates 3--4$\times$ at $\mathrm{DI}\le-90$), indicating a
simulated effective size that is too small; recalibrating on these neutral
markers---rather than on the ancestry-informative markers, which are the test---
lowers simulated differentiation in every DI bin and thus \emph{widens} the
high-DI excess, so the reported magnitudes are conservative.
\end{abstract}

\section{Rationale}

Near-fixation of a parental allele in a hybrid population is expected under a
Dobzhansky--Muller incompatibility or consistent directional selection, but also
under plain drift in small, young populations. The two are separated by
\emph{direction}: selection predicts the same parent fixing repeatedly across
replicates, drift predicts random directions. Quantifying ``more parallel than
drift'' requires a null that reproduces the actual demography, recombination
landscape and sampling design of the data. That is what Module E provides.

\section{Methods}

\subsection{Conditional design}

Rather than simulating the parental species' history, we \emph{condition} on it.
Founders are drawn from real phased parental haplotypes, the mutation rate is set
to zero, and only the post-contact hybrid phase is simulated. This inherits true
allele frequencies, true between-species divergence and---critically---true
within-species haplotype structure, none of which independent binomial draws at
empirical frequencies would reproduce.

\subsection{Founder haplotypes}

Fifteen diploid females per species were imputed and phased \emph{within species}
with Beagle~5.4, informed by the empirical cM map (phasing across the species
boundary would manufacture artificial intermediate haplotypes). Two
\textit{polyctena} individuals with $>$40\% missing data were excluded, giving a
founder pool of 30 \textit{aquilonia} haplotypes and 13 \textit{polyctena} diploid
females over 1{,}114{,}251 markers. Phasing preserved the pool faithfully:
correlation between raw and phased allele frequencies $r=0.999$, expected
heterozygosity $0.216/0.203$ against raw $0.216/0.206$, and mean between-species
$|\Delta p|$ of $0.304$ against $0.303$.

\subsection{Simulation model}

SLiM~5, non-Wright--Fisher, haplodiploid with arrhenotoky: haploid males from
unfertilised eggs, diploid females, spermathecal sperm storage, multiple mating
following Pamilo, and an age-structured life table. The genome is the 26
assembled chromosomes (chromosome 23 excluded), each with its empirical
recombination map. Maps are in cM/Mb and must be converted with
\texttt{scale}=$10^{-8}$ to per-base rates. Mutation rate is zero, so the
simulation redistributes existing real variation rather than generating new
variation. All selection (QTN, DMI, climate) is removed.

Founders are read from one combined per-chromosome VCF containing 30 haploid
\textit{aquilonia} samples and 13 diploid \textit{polyctena} samples. A single
read call per chromosome is required: reading the two species separately creates
two distinct mutations at each position instead of one shared biallelic marker.
Each parameter set is run as 20 independent demes---each an independent founder
draw and independent drift history---sampled every 20 generations to generation
160.

\subsection{Marker sets and DI stratification}

Two marker sets were used. The first is the ancestry-diagnostic set,
$\mathrm{DI}>-25$ (49{,}506 markers). The second, used for the results reported
here, takes the locked Module~A gate (pooled parental
$\mathrm{MAF}\ge0.15$; 650{,}950 markers, DI from $-168$ to $-4$) and draws
\textbf{4{,}000 markers from each of 10 fixed DI bins} (40{,}000 total). Equal
allocation per bin is deliberate: uniform thinning is swamped by the mid-DI bulk
and starves both tails.

This stratification supplies the internal control. Low-DI markers carry little
ancestry information, so ancestry sorting cannot drive their statistics; they
calibrate the demography \emph{in the same runs} that test the high-DI markers.
The test is then the \emph{shape} of the discrepancy across DI: a gap that rises
with DI implicates ancestry-specific processes, whereas a flat gap implicates
demographic misspecification.

\subsection{Requirements for a valid comparison}

Three points proved essential and are worth recording.

\begin{enumerate}\itemsep2pt
\item \textbf{Sample-size matching.} Composite $r^2$ carries a finite-sample bias
of order $1/2n$, and \texttt{sort\_class} likewise fires more readily in small
samples (a 3-diploid sample can near-fix trivially; 50 diploids essentially
cannot). Simulated demes are therefore subsampled to the empirical
per-population sizes (3--20 diploids). Without this the null appears
artificially unsorted and the LD comparison is biased toward younger ages.
\item \textbf{Symmetric treatment.} Any filter applied to the data must be
applied to the null. ``Sorted'' is by definition high among-population
differentiation, so removing sorted loci from the observed data alone
mechanically lowers its $F_{ST}$ and moves it toward the null irrespective of
cause.
\item \textbf{Shared orientation.} Both datasets are polarised with the same
parental reference (the phased founder pool) and the locked Module~A conventions
($\mathrm{MAF}\ge0.15$, $\mathrm{fix\_th}=0.15$, $\mathrm{sort\_th}=0.5$,
$\mathrm{null\_prob}=0.5$).
\end{enumerate}

\section{Controls excluding neutral alternatives}

Before interpreting any excess, several neutral explanations were quantified and
excluded.

\paragraph{Founder heterogeneity.}
If replicate populations were founded from different local source populations,
among-population differentiation would be inflated neutrally. Splitting the
parental pool into $k$ balanced sub-pools and founding each deme from one of
them gives an at-founding among-deme $F_{ST}$ of only $0.0095$--$0.024$
($k=2$--$4$). The absolute ceiling---each deme founded from a single
\textit{aquilonia} haplotype and a single \textit{polyctena} female, the most
extreme founding available---is $0.090$. The mechanism is general: at
ancestry-informative markers both species are near-fixed for opposite alleles,
so any \textit{aquilonia} sub-pool resembles ``aquilonia'' and any
\textit{polyctena} sub-pool ``polyctena''; a 50/50 admixture washes out
within-species structure regardless of how divergent the sources are.

\paragraph{Colony (nest) structure.}
The empirical samples are workers from few colonies (164 workers from 39 nests;
median 2 nests per population; 9 of 20 populations are a single colony), and
nestmates could in principle be siblings, which would depress $\pi$ and inflate
$F_{ST}$. Measured directly, the excess within-nest pairwise genotype
correlation is only $0.047$ ($0.565$ within nests versus $0.517$ between nests in
the same population): nestmates are barely more related than non-nestmates,
as expected for polygynous supercolonial \textit{Formica}. The effect is real but
minor.

\paragraph{Among-population ancestry variance.}
Populations differ in mean ancestry (SD $0.065$, range $0.40$--$0.66$), which
contributes $\mathrm{Var}(a)/\bar{p}(1-\bar{p})\approx0.017$ to $F_{ST}$ at
diagnostic markers---about 11\% of the observed high-DI gap.

\paragraph{Individual outliers.}
Sielva is a near-F$_1$ population (observed heterozygosity $0.813$ against
$0.499$ expected from its ancestry; within-population ancestry SD $0.001$), and
is genuinely young. Removing it changes $\pi$ from $0.299$ to $0.293$ and
$F_{ST}$ from $0.313$ to $0.320$; across all 20 leave-one-out analyses the
largest change in $F_{ST}$ is $0.009$. No single population drives the result.

\section{Results}

\subsection{The neutral null is anchored on markers that cannot sort}

Low-DI markers ($\mathrm{DI}\le-60$) were used to fit the demography by LD decay,
selecting $K=1500$ at generation~160 (Fig.~\ref{fig:ld}). This anchor must be
computed \emph{within} populations and averaged; pooling individuals across
differentiated populations manufactures LD and destroys the comparison.

\subsection{Differentiation scales with DI in the data but not under drift}

Table~\ref{tab:dose} and Fig.~\ref{fig:dose} give the central result. Observed
$F_{ST}$ rises from $0.051$ at the neutral background to $0.315$ at the most
ancestry-informative markers, while the neutral null is nearly flat
($0.191\rightarrow0.156$). The fraction of sorted loci shows the same pattern
more sharply, from a $\sim$1.5-fold excess at low DI to $>$100-fold at high DI.

\begin{table}[htbp]\centering
\caption{Observed and simulated statistics by DI bin at the anchor-matched
neutral cell ($K=1500$, generation 160). ``sorted'' is the fraction of loci
classified as directionally sorted.}
\label{tab:dose}
\small
\begin{tabular}{lrrrrrr}
\toprule
DI bin & $F_{ST}$ obs & $F_{ST}$ null & $\pi$ obs & $\pi$ null & sorted obs & sorted null\\
\midrule
$\le-90$      & 0.051 & 0.191 & 0.448 & 0.381 & 0.026 & 0.021\\
$(-90,-75]$   & 0.257 & 0.185 & 0.308 & 0.344 & 0.184 & 0.103\\
$(-75,-60]$   & 0.303 & 0.183 & 0.260 & 0.314 & 0.350 & 0.233\\
$(-60,-50]$   & 0.307 & 0.179 & 0.251 & 0.313 & 0.398 & 0.267\\
$(-50,-40]$   & 0.309 & 0.172 & 0.253 & 0.328 & 0.378 & 0.183\\
$(-40,-30]$   & 0.305 & 0.167 & 0.277 & 0.364 & 0.260 & 0.072\\
$(-30,-25]$   & 0.312 & 0.161 & 0.293 & 0.388 & 0.175 & 0.030\\
$(-25,-20]$   & 0.310 & 0.160 & 0.297 & 0.398 & 0.150 & 0.015\\
$(-20,-15]$   & 0.314 & 0.161 & 0.302 & 0.405 & 0.130 & 0.001\\
$>-15$        & 0.315 & 0.156 & 0.304 & 0.415 & 0.116 & 0.000\\
\bottomrule
\end{tabular}
\end{table}

\begin{figure}[htbp]\centering
\includegraphics[width=0.85\textwidth]{../Figures/moduleE_fig1_dose_response.pdf}
\caption{Among-population differentiation (top) and the fraction of sorted loci
(bottom, log scale) against DiagnosticIndex bin, for the observed data and the
neutral null. Neutral drift is DI-blind and produces an essentially flat
profile; the data concentrate both differentiation and sorting in
ancestry-informative markers.}
\label{fig:dose}
\end{figure}

\begin{figure}[htbp]\centering
\includegraphics[width=0.68\textwidth]{../Figures/moduleE_fig2_ld_anchor.pdf}
\caption{LD decay on the low-DI anchor markers ($\mathrm{DI}\le-60$), computed
within populations and averaged, for the observed data and the anchor-matched
neutral cell.}
\label{fig:ld}
\end{figure}

\subsection{Why the shape matters}

Drift does not distinguish ancestry-informative from uninformative markers: a
purely demographic error---wrong effective size, wrong age, wrong founding
size---shifts all DI bins together. The simulated profile confirms this,
remaining flat across the full DI range. A discrepancy that \emph{scales} with
ancestry-informativeness therefore cannot be produced by misspecifying the
demography, and the comparison is internally controlled because both datasets
use the same markers, the same MAF gate and the same frequency spectrum.

\section{Limitations and the planned calibration}

\textbf{The null does not yet reproduce the neutral background.} At
$\mathrm{DI}\le-90$ the observed $F_{ST}$ is $0.051$ while the null gives
$0.14$--$0.19$ (depending on $K$): the simulation over-differentiates by
3--4$\times$ on markers where sorting cannot be the cause, and simultaneously
undershoots neutral $\pi$ ($\sim$0.40 against $0.448$). Both signal the same
thing---the \emph{simulated effective size is too small}: isolated demes drift
apart more, and lose more within-population diversity, than the real populations
do at neutral loci. Consequently the \emph{magnitude} of the excess is
provisional, though its \emph{direction} is not: reducing simulated drift lowers
$F_{ST}$ in every DI bin and, because the null is flat across DI, \emph{widens}
the high-DI gap, making the reported values conservative.

\paragraph{Planned analysis.}
The calibration must be performed \emph{on the low-DI markers}---the internal
control that cannot sort---rather than on the ancestry-informative markers (which
are the test) or on LD decay (which, being generated partly by drift, favours the
opposite direction and disagrees with $\pi$/$F_{ST}$). Two demographic parameters
are swept over the full DI-stratified marker set:
\begin{enumerate}\itemsep2pt
\item \textbf{carrying capacity $K$}, extended \emph{upward} beyond the values run
so far (the neutral background requires less drift than even $K=6250$ provides,
so larger effective sizes are needed);
\item \textbf{number of founders} per deme, which sets both within-population
diversity and among-deme founding variance.
\end{enumerate}
The cell that reproduces neutral $F_{ST}$ ($\approx0.051$) and $\pi$
($\approx0.448$) then defines the calibrated null, and the high-DI excess is read
off there. Because the null's $F_{ST}$ is flat across DI, matching the neutral
background at $\sim$0.05 lowers the simulated high-DI $F_{ST}$ to the same level,
so the observed high-DI value ($0.315$) stands against a much lower neutral
expectation than in Table~\ref{tab:dose}; the dose-response is expected to
strengthen, not weaken, under this recalibration. Founding is expected to be
constrained to be \emph{large}: fewer founders lower $\pi$ and raise $F_{ST}$,
moving the neutral background away from the observed values, which itself
excludes a tight-bottleneck neutral explanation on the markers where the test is
trustworthy.

Until the background is reproduced, the honest statement is that
\textbf{differentiation and sorting scale with ancestry-informativeness in a way
neutral drift does not produce}, with the magnitude of that excess to be fixed
once the null reproduces the markers that cannot sort.

\end{document}
